\documentclass[10pt,a4paper]{amsart}
\usepackage[margin=19mm]{geometry}
\usepackage{amsmath,amssymb,mathtools}
\usepackage[scr=rsfs,cal=euler]{mathalfa}
\usepackage{xfrac}
\usepackage[unicode,hidelinks,bookmarks=false]{hyperref}
\newcommand{\ie}{i.e.\ }
\newcommand{\cf}{cf.\ }
\newcommand{\resp}{respectively\ }
\newcommand{\Subsection}{\subsection}
\providecommand{\nobreakdash}{\nobreak\hbox{-}}
\setlength{\emergencystretch}{2em}
\multlinegap0pt
\usepackage{tikz}
\usepackage{xcolor}
\usepackage{amssymb}
\usepackage{algorithm}\usepackage{algpseudocode}
\usepackage{mathtools}
\usepackage[shortlabels]{enumitem}
\newtheorem{lemma}{Lemma}
\newtheorem{proposition}{Proposition}
\usepackage{cleveref}
\crefname{lemma}{Lemma}{Lemmas}
\crefname{proposition}{Proposition}{Propositions}
\DeclareMathOperator{\Pet}{Pet}
\title{Combinatorial proofs of Petrie Pieri rule and Plethystic Pieri rule}
\author{Saintan Wu, Sen-Peng Eu, Kuo-Han Ku, Yu-Sheng Shih}
\date{Source: arXiv:2509.16872v2 (CC BY 4.0)}
\begin{document}
\maketitle

\noindent\emph{Source and licence.} Combinatorial proofs of Petrie Pieri rule and Plethystic Pieri rule. Version arXiv:2509.16872v2 (CC BY 4.0), distributed by arXiv under the Creative Commons Attribution 4.0 International licence, http://creativecommons.org/licenses/by/4.0/, as recorded in the arXiv licence metadata for this exact version and shown on the abstract page https://arxiv.org/abs/2509.16872v2. Full licence notice: \texttt{licenses/arxiv-2509.16872-CC-BY-4.0.txt}.\par\smallskip

\noindent\textit{Editorial scope.} Counted unit: the article's lemma lemma (source0.tex lines 2257--2263). The article's complete original proof of it is delivered with it (source0.tex lines 2264--2364, one \texttt{proof} environment of 101 lines). No mathematics is written, repaired or re-derived: the statement and the proof are the article's own lines, in the article's own notation, and no retained line is substituted or re-flowed.  Retained uncounted as the source-local context the proof needs: lines 1843--1849; lines 2119--2121.  Nothing was omitted from any retained span, so the package carries no omission record.  No retained line contains a citation, so the wrapper carries no bibliography and no bibliography entry.  The retained lines contain the article's own diagram code, which the wrapper typesets with the standard tikz packages; no image file is delivered and the package declares no asset dependency.  Editorial formatting only, in addition: an AMS \texttt{amsart} preamble that restates the article's own declarations and macros used by the retained lines, namely the environments \texttt{lemma}, \texttt{proposition} on separate counters, together with the standard packages those lines need.  The retained environments print in this document as Lemma 1, Proposition 1 in order of appearance, whereas the article prints the same items with the numbering of its own full document.  The complete native source of the article as distributed in the arXiv e-print for this exact version is bundled unchanged as \texttt{sources/185417/source0.tex}.
        \begin{lemma}\label{lemma:cycleCaseLemma}
            Let $A\in \mathcal{CP}(n)$ and suppose that $A$ has row vectors $v_i$. Then 
            \begin{enumerate}
                \item $\sum_i (-1)^i v_i=0$.
                \item The column vectors $v[a'_i,b'_i]$ of $A$ satisfies that $a'_{i+1}+b'_{i+1}=a'_{i}+b'_{i}+2$.
            \end{enumerate}
        \end{lemma}

        \begin{proposition}\label{prop:extendTotallyPetrie}
            Any totally Petrie square matrix can be written in the form $\Pet_k(\lambda,\mu)$.
        \end{proposition}

        \begin{lemma}
            Let $A\in \mathcal{CP}(n)$ and $L(\tau_{\sigma})$ be its corresponding polygonal line. Then
            \begin{enumerate}
            \item The polygonal line $L(\tau)$ does not intersect itself.
            \item If $A$ has size $n\times (n-1)$ and $n$ is even, then $L(\tau)$ is a closed curve, i.e. $a_{n+1}=0$.
            \end{enumerate}
        \end{lemma}

        \begin{proof}
           Notice that the slope of each segment in the polygonal line is positive. For two of the segments to intersect, one's endpoints have coordinates strictly between the other's. This means that in $A$, there is a row whose $1$'s start later but end earlier than another row, which cannot occur in $A$.
        
           When $n$ is even, by \cref{lemma:cycleCaseLemma}, we have an equal number of reversed rows and non-reversed rows in $\sigma$, hence $a_{n+1}=0$.

   % %%Proof not true
   %         For the furthermore part, we make $A$ a square matrix by adding a zero column to the end, and the resulting matrix is $\Pet_k(\lambda,\mu)$ for some $k,\lambda,\mu$ by \cref{prop:extendTotallyPetrie}. The cycle in the $\mathcal{P}'(\Pet_k(\lambda,\mu))$ has even steps and thus is contractible, namely it encloses a polygonal region on the cylinder. 
           
   %          \setcounter{MaxMatrixCols}{20}
   %         \begin{figure}
   %             \centering
   %             \newcommand{\1}{\,\textcolor{red}{1}\,}
   %             $\begin{bmatrix}
   %                 \1&\1&1& \\
   %                 \1&\1&1&1\\
   %                   &  &1&1&\1&1\\
   %                   &  & & & &1&1&\1&\1& \\
   %                   &  & & & & &1&\1&\1&1&\\
   %                   &  & & & & & &\1&\1&1&1&\\
   %                   &  & & & & & & & & &1&1&\1&\1& \\
   %                   &  & & & & & & & & & &1&\1&\1&1\\
   %                   &  & & & & & & & & & & &\1&\1&1&\1\\
   %             \end{bmatrix}$
   %         \begin{tikzpicture}[baseline=-.5cm, scale=0.8]
   %              \filldraw[olive!60] (1.5,1.4)--(1.8,1.7)-- (2.2,1.7)--(2.2, 1)-- (1.8,1)-- (1.5,.7)--cycle;
   %              \filldraw[olive!60] (3.8, .9) -- (5, .9)--(5.3,1.2)--(5.3,0.3)--(5,0)--(4.6,0)--(4.6,-1)--(5,-1)--(5.3,-.7)--(5.3,-1.7)--(5,-2)--(3.8,-2)--cycle;
   %              \filldraw[olive!60] (1.5,-.3)--(1.8,0)--(3,0)--(3,-2)--(2.2,-2)--(2.2,-1)--(1.8,-1)-- (1.5,-1.3)--cycle;
			% 	\draw [red] (1.5,1.4)--(1.8,1.7)-- (4.5,1.7) ;
   %              \draw [red] (3.7, .9) -- (5, .9)--(5.3,1.2);
			% 	\node at (3.4,1.6) [above]{$\leq -6$};
				
   %              \draw (1.8,-2.7)
   %              %--(2.1,-2.7)--(2.1,-1.7)--(2.3,-1.7)--(2.3,-2.7)
   %              --(5, -2.7);
			% 	\node at (3.4,-2.6) [below]{$\geq 9$};
				
			% 	% \draw (3.8, 1) circle (6pt) node[inner sep=2pt] (-8){\scriptsize -8};
			% 	\filldraw[fill=white] (3, 1) circle (6pt) node[inner sep=2pt] (-7){\scriptsize -5};
			% 	\filldraw[fill=white] (2.2, 1) circle (6pt) node[inner sep=2pt] (-6){\scriptsize -4};
			% 	\filldraw[fill=white] (4.6, 0) circle (6pt) node[inner sep=2pt] (-4){\scriptsize -3};
			% 	\filldraw[fill=white] (3.8, 0) circle (6pt) node[inner sep=2pt] (-3){\scriptsize -2};
			% 	\filldraw[fill=white] (3, 0) circle (6pt) node[inner sep=2pt] (-2){\scriptsize -1};
			% 	\filldraw[fill=white] (2.2, 0) circle (6pt) node[inner sep=2pt] (-1){$0$};
			% 	\filldraw[fill=white] (4.6,-1) circle (6pt) node[inner sep=2pt] (1){$1$};
			% 	\filldraw[fill=white] (3.8,-1) circle (6pt) node[inner sep=2pt] (2){$2$};
			% 	\filldraw[fill=white] (3  ,-1) circle (6pt) node[inner sep=2pt] (3){$3$};
			% 	\filldraw[fill=white] (2.2,-1) circle (6pt) node[inner sep=2pt] (4){$4$};
			% 	\filldraw[fill=white] (4.6,-2) circle (6pt) node[inner sep=2pt] (6){$5$};
			% 	\filldraw[fill=white] (3.8,-2) circle (6pt) node[inner sep=2pt] (7){$6$};
			% 	\filldraw[fill=white] (3  ,-2) circle (6pt) node[inner sep=2pt] (8){$7$};
			% 	\filldraw[fill=white] (2.2,-2) circle (6pt) node[inner sep=2pt] (9){$8$};
			% 	\draw[red,-] (-6) -- (1.8,1)-- (1.5,.7);\draw[red,-] (5.3,.3)--(5.,0)--(-4);
			% 	\draw[red,-] (-2) -- (-1);
			% 	\draw[red,-] (-1) -- (1.8,0)-- (1.5,-.3);\draw[red,-] (5.3,-.7)--(5.,-1)--(1);
			% 	\draw[red,-] (4) -- (1.8,-1)-- (1.5,-1.3);\draw[red,-] (5.3,-1.7)--(5.,-2)--(6);
			% 	\draw[red,-] (6) -- (7);
			% 	\draw[red,-] (8) -- (9);
                
			% 	\draw[->] (3, 1.7) -- (-7);
			% 	\draw[->] (2.2, 1.7) -- (-6);
			% 	\draw[<-] (3.8,.9) -- (-3);
   %              \draw[->] (-4) -- (1);
			% 	\draw[<-] (-3) -- (2);
   %              \draw[->] (-2) -- (3);
			% 	\draw[<-] (2) -- (7);
			% 	\draw[->] (3) -- (8);
			% 	\draw[<-] (4) -- (9);
			% \end{tikzpicture}
   %          \begin{tikzpicture}[scale=0.5, baseline=2.25cm]
   %                  \draw[fill=blue!20]
   %                      (0,0)
   %                      --(2,3) 
   %                      --(6,6) 
   %                      --(1,3) 
   %                      --(5,6) 
   %                      --(10,9) 
   %                      --(7,6) 
   %                      --(11,9) 
   %                      --(9,6) 
   %                      --(4,3) 
   %                      --(8,6) 
   %                      --(3,3)  
   %                      --cycle;
   %                     \fill (0,0) circle (3pt);
   %                     \fill (2,3) circle (3pt);
   %                     \fill (6,6) circle (3pt);
   %                     \fill (1,3) circle (3pt);
   %                     \fill (5,6) circle (3pt);
   %                     \fill (10,9) circle (3pt);
   %                     \fill (7,6) circle (3pt);
   %                     \fill (11,9) circle (3pt);
   %                     \fill (9,6) circle (3pt);
   %                     \fill (4,3) circle (3pt);
   %                     \fill (8,6) circle (3pt);
   %                     \fill (3,3) circle (3pt);
   %              \end{tikzpicture}
   %             \caption{Caption}
   %             \label{fig:placeholder}
   %         \end{figure}
            
        \end{proof}

\end{document}
